\chapter*{Introduction}
\addcontentsline{toc}{chapter}{Introduction}
\label{ch:introduction}

The Item~1A risk-factor section of the U.S.\ 10-K, mandatory since
2005, is the longest and most carefully written piece of
risk-disclosure language in U.S.\ corporate reporting. Two parallel
literatures argue that this section conveys economically meaningful
information about firm-level risk, but they reach opposite
predictions about its sign. \textcite{campbell2014} show that
\emph{larger} risk-factor sections, both in raw word count and in
the count of risk-related words, are associated with higher
subsequent return volatility, market beta, and bid--ask spreads:
when managers write more about risk, they are signalling more risk.
\textcite{hope2016} show that, for a given length, risk-factor
sections that contain a higher proportion of \emph{specific}
content---names, numbers, geographies---are associated with
stronger analyst-forecast revisions and lower implied cost of
capital: more concrete risk language pins down exposures and reduces
uncertainty.

Much of the empirical literature measures both predictions with a
single statistic: the dictionary-based risk-word density of
Item~1A. This is the share of the section's words that match the
\textcite{loughran2011} uncertainty and litigious word lists.
But the ratio has section length in its denominator. The volume
channel sits in the numerator and the specificity channel sits in
the denominator, so the two effects are mixed into one number that
has been hard to sign cleanly in the data.

This thesis takes a small methodological step with a real payoff.
It splits the dictionary risk-density into its two parts---the
natural log of the Item~1A word count, and the risk-word density
itself---and enters them jointly in a panel regression of
post-filing log-volatility. With both regressors in the equation,
length captures the volume channel net of any change in boilerplate
share, and density captures the specificity channel net of any
change in volume. Both predictions can therefore be tested in the
same regression, and the data can choose between them.

On an unbalanced panel of 535 large U.S.\ non-financial firms over
fiscal years 2010--2024, the joint specification at the 30-day
horizon produces both predictions: longer Item~1A sections predict
higher post-filing volatility (the volume effect, H1), while
denser Item~1A sections, holding length constant, predict lower
volatility (the specificity effect, H2). Both signs and
significance hold at every post-filing horizon between 5 and
365~days and on both pre-COVID and COVID-era sub-samples. The
joint specification also shows that the dictionary-only estimate
of the specificity channel is roughly twice as large as the true
effect; the omitted volume channel loads onto density through the
negative correlation between length and density.

A second-stage analysis looks inside the LM risk lexicon itself.
The composite RiskDensity used for H2 sums two LM word lists with
very different production technologies. Splitting it into
LM-Uncertainty and LM-Litigious sub-densities shows that the whole
negative effect comes from the LM-Litigious sub-list. At the
30-day horizon, holding length fixed, LitDensity has a coefficient
of $-6.16$ ($t = -3.38$) on post-filing log-volatility, while
UncDensity is small and not statistically different from zero. The
pattern holds at every horizon and on both COVID sub-samples. The
reading is the one suggested by
\textcite{nelson2007,cazier2021,beatty2019}: litigious-list
language is standardised cautionary boilerplate produced under the
safe-harbour regime of the 1995 Private Securities Litigation
Reform Act, and is typical of mature, legally supervised firms.
The decomposition is treated as a refinement of H2 rather than as
a separate hypothesis: it shows which part of the LM lexicon
drives the H2 specificity effect.

A third hypothesis (H3) introduces a peer-relative measure of
strategic silence motivated by the classical disclosure-equilibrium
literature \parencite{verrecchia1983,dye1985}. For each firm-year
we build an \emph{omission gap}: the share of TF--IDF disclosure
mass written by the firm's two-digit-SIC industry peers in the
same fiscal year that the focal firm itself does not mention
(Section~\ref{ssec:silence}). Holding Item~1A length, RiskDensity,
and the H1 and H2 controls fixed, a higher omission gap is
positively associated with post-filing log-volatility under firm
and year fixed effects ($\hat{\beta}^{\text{om}} = +1.345$,
$t = +2.93$ at the 30-day horizon). The result survives an
industry-shuffle placebo, gets stronger with stricter length
controls and a large-peer-cell restriction, and remains significant
when the omission gap is lagged by one fiscal year. To our
knowledge, the H3 evidence is the first panel test of strategic
silence in the post-filing realised-volatility literature, with
two qualifications discussed in Chapter~\ref{ch:robustness}.

The remainder of the thesis is organised as follows.
Chapter~\ref{ch:literature} reviews the volume, specificity,
boilerplate, and strategic-silence literatures and motivates the
three hypotheses.
Chapter~\ref{ch:data} describes the SEC EDGAR data, the textual
extraction pipeline, and the construction of the financial controls
and the post-filing volatility windows. Chapter~\ref{ch:methodology}
formalises the joint specification, the LM-list decomposition, and
the peer-relative omission specification.
Chapter~\ref{ch:results} reports the headline tests of H1, H2, and
H3. Chapter~\ref{ch:robustness} reports the multi-horizon
evidence, pre/post-COVID splits, a within-firm fixed-effects
specification, the H3 robustness battery (placebo,
length-residualised, large-cell, decile-rank, and lead-lag
specifications), and a transparent account of seven alternative
specificity-channel candidates that were tested and rejected. The
Conclusion summarises the contribution and discusses limitations
and suggested extensions.
